\subsection{Simulated populations, observed host forcing, and experimental design}
\label{num-esm-simworld}

The available data do not isolate the effect of host variability at a fixed
mean on parasite persistence. We therefore use controlled host schedules to
separate variability from residence in a steady stage, and observed host
sequences to examine where those predictions cease to apply. The numerical
results concern the specified parasite models. Host data determine forcing
and checks of the predictions, without fitting parasite parameters or adding
a mechanism to reproduce an observed pattern.

\subsubsection{Model and measured outcomes}
\label{num-esm-model}
% Sources: virtualworld/protocol.md, including both dated fix1 addenda; report.md.

The simulated world represents individuals in ordered biological stages.
Conditional on the environment, individuals undergo independent demographic
events. The pilot uses two geometric residence distributions (a constant
exit probability at each step), an exposed mean stay of 2 steps, and a steady
mean stay $s\in\{1,5,20\}$ steps. An individual leaving the exposed stage
produces a Poisson number of entrants to the steady stage, with mean
$h_t=H_t/\bar H$ (host abundance relative to its stationary arithmetic mean).
The return transfer has mean offspring yield $R_0$, the constant-environment
reproduction number. Scarcity reduces successful output without postponing
departure. There is no additional stage mortality, host feedback, saturation,
alternative host, dispersal, or density regulation in the pilot.

Host forcing is generated as
\begin{equation}
 h_t=\exp(\sigma z_t-\sigma^2/2),\qquad
 z_t=\phi z_{t-1}+\sqrt{1-\phi^2}\,\varepsilon_t,
 \label{num-host-forcing}
\end{equation}
where $z_t$ has its stationary standard normal distribution and
$\varepsilon_t$ are independent standard normal innovations. Thus $\sigma$
is the stationary standard deviation of log abundance and $\phi$ is its
lag-one correlation. The stationary arithmetic mean of $h_t$ is one.
Finite paths are not rescaled, and their realized means and log variances
are retained. This distinction matters when comparing an imposed schedule
with a finite experimental exposure.

Extinction is the first step at which the total population reaches zero.
All populations, including extinct ones, contribute to extinction estimates.
Early growth is calculated from the conditional-mean population operator
over steps 0--100 for every host schedule. It is not individual growth
conditioned on survival. The threshold loss $D$ is the difference between
the log offspring multipliers that restore zero long-run invasion growth
under fluctuating hosts and under the arithmetic-mean-rate operator.
The multiplier is applied only at circuit completion. We solve for this
root separately from finite-horizon extinction and growth.

\subsubsection{Validation before the pilot}
\label{num-esm-validation}
% Sources: virtualworld/protocol.md, validation.md and its correction, report.md;
% verification/replication/validation-rerun.md.

All five of five registered gates passed before the pilot. The gates link
individual demographic draws to the conditional-mean operator, then compare
that operator with independently calculated limits and earlier numerical
results (Table~\ref{num-validation-table}). A second party wrote separate
code from the protocol, without importing or copying the original
implementation, and repeated every gate with independent random seeds.
Both implementations passed 190 of 190 comparisons. Agreement concerns
these laws and tolerances, rather than the biological realism of every
available model extension.

\begin{table}[htbp]
\centering
\caption{Registered validation gates and independent repetition. Counts include
algebraic, precision, root, and directional checks within each gate.}
\label{num-validation-table}
\begin{tabular}{p{0.57\linewidth}rr}
\toprule
Scope & Original & Independent \\
\midrule
Small-noise ring coefficient, band, and individual means & 53/53 & 53/53 \\
Mechanism extensions and conditional-mean matrices & 30/30 & 30/30 \\
Exact gamma-host case and adjacent-stay decreases & 54/54 & 54/54 \\
Positive environmental memory and adjacent decreases & 44/44 & 44/44 \\
Constant-host demographic extinction & 9/9 & 9/9 \\
\midrule
Total & 190/190 & 190/190 \\
\bottomrule
\end{tabular}
\end{table}

The first gate used stage means $(2,5)$, $(1,2,7)$, and $(2,3,4,5)$ at
$\phi=0$ and $0.6$, with $\sigma=0.03$. Algebraic coefficient error was
required to be at most $10^{-10}$. Simulated coefficients used 32 independent
antithetic pairs (opposite Gaussian innovations), 65,536 recorded steps,
and 2,048 burn-in steps. The allowed error was $0.05K+4\,\mathrm{SE}$,
with $\mathrm{SE}\leq0.03K$, where $K$ is the second-order coefficient in
$D\simeq\sigma^2K$ and $\mathrm{SE}$ denotes standard error.

The second gate checked six of six selected mechanism settings, covering
the base model, residence shape, alternative hosts, patches, mortality,
and delaying contact. Matrix agreement was required within $10^{-11}$.
Threshold agreement allowed four combined standard errors plus $0.002$,
with separate precision, root, and half-run checks. The third gate used
the exact gamma-host model of \citet{roerdink1988}. Gamma shapes
$2,1,0.25$ were crossed with staying probabilities $0.2,0.5,0.8,0.95$.
Independent quadrature refinement had to change roots by at most $10^{-8}$,
and simulation had to agree within four standard errors plus $0.0005$.
All nine of nine adjacent exact and simulated decreases passed their
respective directional checks.

The fourth gate crossed $\sigma=0.3,0.6$, $\phi=0.6,0.9$, and steady
means $2,5,10$, with exposed mean one. All twelve of twelve roots agreed
with the archived references within four combined standard errors plus
$0.001$, and eight of eight adjacent decreases exceeded four paired
standard errors. These are finite-grid checks. The fifth gate compared
10,000 trajectories per reproduction setting with a separately implemented
probability-generating-function recursion (an exact recursion for the
specified branching probabilities). The settings were $\ln R_0=-0.1,0,0.1$,
initial counts $(10,0)$, stays $(2,5)$, and endpoints 20, 60, and 120 steps.
All nine of nine comparisons met the tolerance
$4\sqrt{p(1-p)/10000}+0.002$, where $p$ is the reference extinction probability.

Sampling intervals do not bound all systematic error. The dated validation
correction identifies directional statistics as differences of corrected
influence-sample means, rather than differences of the printed uncorrected
roots. The registered root tolerance was $10^{-7}$, and the correction
changed no gate decision. The independent report also distinguishes passing
registered tolerances from falling inside another run's pointwise interval.

\subsubsection{Pilot and corrected initialization}
\label{num-esm-pilot}
% Sources: virtualworld/pilot.md and dated corrections; pilot-rerun-fix1.md;
% fix1-report.md; verification/replication/pilot-rerun.md.

The factorial pilot crossed $\sigma=0,0.3,0.6$, $\phi=0,0.6$, and
$s=1,5,20$. Each of 18 cells contained 200 independent host and demographic
replicates, giving 3,600 of 3,600 completed trajectories at 500 steps.
Initial abundance was 20 and $\ln R_0=0.05$. The two correlation labels
at zero noise are separate controls with independent demographic draws.
They were not pooled.

The original integer allocation assigned remainders to the most abundant
stage. The corrected run assigns them to the largest fractional remainders,
giving initial stage counts $(13,7)$, $(6,14)$, and $(2,18)$ for the three
steady stays. The same seed schedules were reused. Consequently, the
original and corrected runs are not independent, and their difference
cannot be attributed solely to allocation rounding. The repair also removed
an artificial population ceiling, using sums of independent count draws,
and clarified movement, finite-record centering, and a degenerate
one-step residence distribution. All 22 of 22 final tests and 144 of 144
saved-artifact checks passed. The original external-reference gates were
retained, while the changed sampler was checked again against the known
branching law.

Corrected extinction and growth estimates are given in
Tables~\ref{num-pilot-extinction} and \ref{num-pilot-growth}.
Figure~\ref{num-1} deliberately uses the original archived cell file.
Its nine plotted cells retain their original counts and intervals, so
the figure must not be used to reconstruct the corrected contrast.

\begin{table}[htbp]
\centering
\caption{Corrected pilot extinction at 500 steps. Each row contains 200 independent
populations. Intervals are pointwise 95\% score intervals.}
\label{num-pilot-extinction}
\begin{tabular}{rrrrl}
\toprule
$\sigma$ & $\phi$ & $s$ (steps) & Extinct/total & Probability [interval] \\
\midrule
0 & 0 & 1 & 71/200 & $0.355\ [0.292, 0.423]$ \\
0 & 0 & 5 & 71/200 & $0.355\ [0.292, 0.423]$ \\
0 & 0 & 20 & 44/200 & $0.220\ [0.168, 0.282]$ \\
0 & 0.6 & 1 & 68/200 & $0.340\ [0.278, 0.408]$ \\
0 & 0.6 & 5 & 73/200 & $0.365\ [0.301, 0.434]$ \\
0 & 0.6 & 20 & 47/200 & $0.235\ [0.182, 0.298]$ \\
0.3 & 0 & 1 & 104/200 & $0.520\ [0.451, 0.588]$ \\
0.3 & 0 & 5 & 79/200 & $0.395\ [0.330, 0.464]$ \\
0.3 & 0 & 20 & 48/200 & $0.240\ [0.186, 0.304]$ \\
0.3 & 0.6 & 1 & 144/200 & $0.720\ [0.654, 0.778]$ \\
0.3 & 0.6 & 5 & 103/200 & $0.515\ [0.446, 0.583]$ \\
0.3 & 0.6 & 20 & 51/200 & $0.255\ [0.200, 0.320]$ \\
0.6 & 0 & 1 & 182/200 & $0.910\ [0.862, 0.942]$ \\
0.6 & 0 & 5 & 125/200 & $0.625\ [0.556, 0.689]$ \\
0.6 & 0 & 20 & 68/200 & $0.340\ [0.278, 0.408]$ \\
0.6 & 0.6 & 1 & 194/200 & $0.970\ [0.936, 0.986]$ \\
0.6 & 0.6 & 5 & 157/200 & $0.785\ [0.723, 0.836]$ \\
0.6 & 0.6 & 20 & 86/200 & $0.430\ [0.363, 0.499]$ \\
\bottomrule
\end{tabular}
\end{table}

\begin{table}[htbp]
\centering
\small
\setlength{\tabcolsep}{3pt}
\caption{Corrected pilot conditional-mean early growth per step (200 schedules
per cell), threshold loss $D$ (32 independent antithetic pairs per cell),
and the second-order prediction $\sigma^2K$. Brackets give pointwise 95\%
$t$ intervals for growth and $D$. The dated root correction does not affect
the displayed threshold precision.}
\label{num-pilot-growth}
\begin{tabular}{rrrlll}
\toprule
$\sigma$ & $\phi$ & $s$ & Early growth [interval] & $D$ [interval] & $\sigma^2K$ \\
\midrule
0 & 0 & 1 & $0.016766\ [0.016766, 0.016766]$ & $0.000000\ [0.000000, 0.000000]$ & $0.000000$ \\
0 & 0 & 5 & $0.007222\ [0.007222, 0.007222]$ & $0.000000\ [0.000000, 0.000000]$ & $0.000000$ \\
0 & 0 & 20 & $0.002318\ [0.002318, 0.002318]$ & $0.000000\ [0.000000, 0.000000]$ & $0.000000$ \\
0 & 0.6 & 1 & $0.016766\ [0.016766, 0.016766]$ & $0.000000\ [0.000000, 0.000000]$ & $0.000000$ \\
0 & 0.6 & 5 & $0.007222\ [0.007222, 0.007222]$ & $0.000000\ [0.000000, 0.000000]$ & $0.000000$ \\
0 & 0.6 & 20 & $0.002318\ [0.002318, 0.002318]$ & $0.000000\ [0.000000, 0.000000]$ & $0.000000$ \\
0.3 & 0 & 1 & $0.011250\ [0.009876, 0.012623]$ & $0.015410\ [0.015332, 0.015488]$ & $0.015000$ \\
0.3 & 0 & 5 & $0.006235\ [0.005595, 0.006874]$ & $0.006642\ [0.006608, 0.006676]$ & $0.006429$ \\
0.3 & 0 & 20 & $0.002083\ [0.001897, 0.002270]$ & $0.002123\ [0.002114, 0.002132]$ & $0.002045$ \\
0.3 & 0.6 & 1 & $0.007101\ [0.004275, 0.009926]$ & $0.028818\ [0.028686, 0.028950]$ & $0.028846$ \\
0.3 & 0.6 & 5 & $0.005696\ [0.004480, 0.006913]$ & $0.015984\ [0.015910, 0.016058]$ & $0.015836$ \\
0.3 & 0.6 & 20 & $0.002128\ [0.001722, 0.002533]$ & $0.005508\ [0.005472, 0.005544]$ & $0.005408$ \\
0.6 & 0 & 1 & $-0.004692\ [-0.007661, -0.001723]$ & $0.066300\ [0.066034, 0.066567]$ & $0.060000$ \\
0.6 & 0 & 5 & $0.002464\ [0.001326, 0.003602]$ & $0.029146\ [0.029015, 0.029277]$ & $0.025714$ \\
0.6 & 0 & 20 & $0.001812\ [0.001342, 0.002282]$ & $0.009602\ [0.009552, 0.009651]$ & $0.008182$ \\
0.6 & 0.6 & 1 & $-0.021369\ [-0.027088, -0.015649]$ & $0.115641\ [0.115115, 0.116167]$ & $0.115385$ \\
0.6 & 0.6 & 5 & $-0.002202\ [-0.004551, 0.000146]$ & $0.064173\ [0.063701, 0.064644]$ & $0.063345$ \\
0.6 & 0.6 & 20 & $0.000766\ [-0.000038, 0.001569]$ & $0.023372\ [0.023223, 0.023522]$ & $0.021631$ \\
\bottomrule
\end{tabular}
\end{table}


For the four independent cells at $\phi=0.6$, define the interaction
\begin{equation}
 \Delta=\{P_{\mathrm{ext}}(0.6,1)-P_{\mathrm{ext}}(0,1)\}
       -\{P_{\mathrm{ext}}(0.6,20)-P_{\mathrm{ext}}(0,20)\},
 \label{num-interaction}
\end{equation}
where $P_{\mathrm{ext}}(\sigma,s)$ is extinction probability at the declared
horizon. Positive $\Delta$ indicates a smaller noise-associated extinction
increment at the longer stay. Corrected counts in that equation are
$194/200$, $68/200$, $86/200$, and $47/200$, respectively. They give
$\Delta=0.435$ (95\% normal interval $[0.320852,0.549148]$).
Independent beta-posterior draws with one-half added to both binomial
counts give a sensitivity interval $[0.31869,0.54693]$. This is a posterior
sensitivity calculation, not another frequentist confidence interval.
The analogous early-growth loss interaction is $0.036582$ per step
($[0.0308415,0.0423226]$). The corrected calendar-time planning estimate
is 29 populations per cell, but it does not set the primary design.

The full independent rerun completed 18 of 18 cells with 400 populations
per cell, or 7,200 trajectories. Its interaction was $0.425$
($[0.346760,0.503240]$). All 18 of 18 extinction estimates fell inside
the corrected pilot's score intervals. This rerun and the separate audit
below used independent seeds and must not be treated as the same sample.
Thresholds used 32 independent antithetic pairs per cell, with 32,768
recorded steps and 2,048 burn-in steps. All 18 of 18 corrected-pilot
threshold estimates matched the original values because integer
initialization does not enter that calculation. Their intervals describe
sampling across pairs and exclude systematic finite-time error.

\subsubsection{Matching generations and planning replication}
\label{num-esm-matched}
% Sources: verification/vworld-check/vworld-check.md; virtualworld/matched-generations-fix1.md;
% fix1-report.md and protocol.md dated completion addendum.

At 500 steps, stays one and 20 represent approximately 166.7 and 22.7
reference generations. Subtracting constant-host extinction removes the
constant-host truncation difference, but leaves unequal exposure to noisy
generations. We define reference generation time as $T_{\mathrm{gen}}=2+s$
steps and compare at a declared number $G$ of those generations.
This is not the realized lineage-generation count under noise.

With 400 populations per cell, the independent audit gave an interaction
of $0.4075$ at 500 steps (95\% interval $[0.3291,0.4859]$).
Its matched-generation estimates ranged from $0.220$ to $0.298$, with
five of five pointwise intervals above zero (Table~\ref{num-audit-matched}).
The approximately 22.7-generation row is a slice of the calendar-time run.
The remaining horizons share trajectories from a fresh run, so estimates
at different horizons are correlated.

\begin{table}[htbp]
\centering
\caption{Independent audit at matched reference generations, with 400 populations
per cell in the four-cell comparison. Intervals are pointwise 95\% normal intervals.}
\label{num-audit-matched}
\begin{tabular}{rrrl}
\toprule
$G$ & Steps, $s=1$ & Steps, $s=20$ & Interaction [interval] \\
\midrule
$\simeq22.7$ & 68 & 500 & $0.238\ [0.151,0.324]$ \\
23 & 69 & 506 & $0.220\ [0.131,0.309]$ \\
50 & 150 & 1100 & $0.288\ [0.198,0.377]$ \\
100 & 300 & 2200 & $0.298\ [0.213,0.382]$ \\
167 & 501 & 3674 & $0.243\ [0.160,0.325]$ \\
\bottomrule
\end{tabular}
\end{table}

The same audit examined $\ln R_0=0.2$ in four fresh cells with 400
populations each. At 500 calendar steps, extinction counts were $6/400$
and $223/400$ for constant and noisy hosts at stay one, and $5/400$ and
$25/400$ at stay 20. The interaction was $0.4925$
($[0.4360,0.5490]$). Both stays had positive long-run invasion growth
under this setting. This comparison establishes dependence on the
reproduction regime, without mapping a full response curve or establishing
a monotone relation with $R_0$.

The subsequent corrected matched-generation run used 400 populations per
cell, or 1,600 trajectories, followed through $G=100$.
The single confirmatory horizon $G=50$ was registered before running.
The other horizons describe the time course with pointwise intervals
(Table~\ref{num-planning}). For independent cells and equal allocation,
the registered normal approximation was
\begin{equation}
 n_{80}=\left\lceil
 \frac{(z_{0.975}+z_{0.80})^2\sum_{i=1}^{4}p_i(1-p_i)}{\Delta^2}
 \right\rceil,
 \label{num-power}
\end{equation}
where $p_i$ is the extinction probability in cell $i$, $z_q$ is the
standard normal quantile of probability $q$, and the target is 80\% power
at two-sided significance level $0.05$.

\begin{table}[htbp]
\centering
\caption{Corrected matched-generation run and model-conditional planning.
All four cells contain 400 populations. The final column is the 10th--90th
percentile range from beta-parameter sensitivity, not an interval for achieved power.}
\label{num-planning}
\begin{tabular}{rllrr}
\toprule
$G$ & Step horizons & Interaction [95\% interval] & $n_{80}$/cell & Range/cell \\
\midrule
23 & 69 / 506 & $0.3325\ [0.2446,0.4204]$ & 58 & 41--85 \\
50 & 150 / 1100 & $0.3275\ [0.2383,0.4167]$ & 61 & 44--91 \\
100 & 300 / 2200 & $0.3425\ [0.2571,0.4279]$ & 51 & 38--73 \\
\bottomrule
\end{tabular}
\end{table}

The maximum estimate is 61 populations in each of four cells, or 244
populations. It concerns marginal power at each horizon under these
estimated probabilities. It does not establish joint power, familywise
error control, or the sample size of a biological experiment. Observation
error, biological blocks, and calibration uncertainty remain to be included.
The beta-parameter calculation used 100,000 interaction draws and 10,000
draws for planning sensitivity. All three of three interaction sensitivity
intervals excluded zero.

The first matched-run report incorrectly treated a right-censoring sentinel
at step 301 as an extinction by the secondary endpoint of 500 steps.
Its secondary row was withdrawn. The three of three primary matched-horizon
rows remained valid. The affected seeded paths were subsequently completed
through 500 steps after a horizon guard and regression check were added.
Earlier host values and population totals matched exactly. The replacement
secondary interaction was $0.4875$ ($[0.4073,0.5677]$), with a conditional
planning estimate of 23 per cell. It is not used for primary planning,
and the replay adds no independent replicates.

\subsubsection{Observed host trajectories and departures}
\label{num-esm-host-series}
% Sources: virtualworld/q4/report.md and final audit; verification/q4-check/q4-check.md.

The host-series campaign selected 120 of 1,331 eligible compilation records
and included 71 valid \textit{Tribolium} stage records. These comprise 191
record labels but 183 distinct abundance sequences, because eight control
records are shared between experiments. Of the 191 labels, 127 contain
only positive observations and 64 contain zeros. One adult record was
excluded after input validation detected an impossible negative count.
Zeros in valid records were retained without pseudocounts or interpolation.

Each record was normalized by its arithmetic mean and replayed circularly.
One model step equals one observation interval, with annual and subannual
records kept in separate sampling strata. Stage means were $(2,s)$ or
$(2,2,s)$, with $s=1,5,20$. Thus 191 records produced 1,146 cells.
The threshold is calculated from the dominant eigenvalue of the matrix
product over the repeated record. It is independent of starting phase.
This construction does not estimate the long-run law of the natural
process that produced a finite record. Biological timescales are not
inferred from its observation interval.

The independent implementation reproduced 204 of 204 cells in a stratified
subsample and 1,146 of 1,146 cells in the full audit. Maximum absolute
threshold disagreement was $8.01\times10^{-13}$, below the registered
absolute criterion $10^{-4}$ (an alternative criterion allowed relative
error below 1\%). All 144 of 144 audited branching-probability endpoints
also agreed, with maximum error $1.68\times10^{-14}$.

We compared $D$ with a second-order prediction based on first-order
autoregressive memory and with one using all circular log autocovariances.
The tolerance was the larger of $0.01$ and 20\% of the absolute predicted
value. Among 762 finite-log cells, 224 of
762 departed from the autoregressive prediction and 211 of 762 from the
full-covariance prediction. The other 384 of 1,146 cells contain zeros
and have no finite log variance. They cannot count as agreement with a
log-host expansion.

Longer stays reduced $D$ in 709 of 764 adjacent comparisons, with 55 of
764 reversals and zero of 764 flat comparisons. For two stages, decreases
occurred in 158 of 191 comparisons from stay one to five and 172 of 191
from five to 20. For three stages, the respective counts were 188 of 191
and 191 of 191. All 55 of 55 reversals disappeared in each of eight
random permutations of the corresponding host record. Permutation retains
the abundance distribution while changing temporal order.

The independent primary ledger contains 271 departures
(Table~\ref{num-departures}). It separates a need to extend the assumptions
from a departure whose cause remains unresolved. All four of four
positive-record reversals have the direction predicted by the existing
spectral response. Lengthening a stay can remove a fluctuation benefit
associated with host timing, raising a negative $D$ toward zero.
The 51 zero-containing reversals are in larval records, including seven
shared-control comparisons, leaving 44 distinct-sequence reversals.

\begin{table}[htbp]
\centering
\caption{Primary scientific departure ledger. Category counts use the full
ledger of 271 departures as denominator, rather than independent biological replicates.}
\label{num-departures}
\begin{tabular}{p{0.57\linewidth}rl}
\toprule
Departure class & Count & Classification \\
\midrule
Positive-record stay reversal & 4/271 & Existing results \\
Zero-containing stay reversal & 51/271 & A13 extension \\
Full-covariance discrepancy, phase-insensitive & 21/271 & Existing results \\
Full-covariance discrepancy, phase-sensitive & 190/271 & Existing results \\
Fluctuating-versus-constant treatment discrepancy & 5/271 & Existing results \\
\midrule
Explained by existing results & 220/271 & \\
Requiring the hosts-with-zeros extension & 51/271 & \\
Unresolved & 0/271 & \\
\bottomrule
\end{tabular}
\end{table}

The initial report left 94 of 211 full-covariance departures unresolved
after a one-sided amplitude reduction to $0.1$. The independent check
used symmetric amplitude reductions to cancel odd-order contributions.
All 211 of 211 passed at symmetric amplitude $0.01$, and all 211 of 211
passed the one-sided check at $0.001$ under the 5\% diagnostic tolerance.
The remaining difference at the original amplitude is therefore compatible
with the limits of a second-order approximation. Phase randomization
showed sensitivity beyond covariance in 190 of 211 cells, but it does not
uniquely separate marginal shape from higher-order temporal dependence.

For fluctuating-versus-constant flour treatments, both predictors gave the
correct direction in 36 of 36 evaluable adult contrasts. Magnitude agreement
held in 31 of 36, with five of 36 departures explained by finite amplitude
and effects beyond covariance. Another 36 of 72 total contrasts involved
zero-containing larval series and were inapplicable to log-based predictions.
Six of 36 evaluable contrasts used two fluctuating cultures against three
controls after the invalid-record exclusion. Restricting to complete-source
contrasts gives 25 of 30 agreements and five of 30 departures.

The original departure file contains 501 overlapping entries. In addition
to the 271 primary entries, it includes 224 autoregressive prediction
departures, five duplicate treatment departures under that predictor, and
one demographic validation discrepancy. That last comparison remains a
registered failure: one of 64 simulated populations survived where the
predicted probability was $0.000112655781$. An exact calculation gave
probability $0.00718451$ of at least one survivor. This traces the
discrepancy to a discrete rare-event diagnostic without changing the
original 479-of-480 pass count. Integer simulation had also stalled on
large populations. Its incomplete paths were not classified, and an
amended, separately validated generating-function calculation completed
persistence outcomes under the same branching law. No tolerance was relaxed.

\subsubsection{Hosts with zeros}
\label{num-esm-zeros}
% Source: a13/report.md overrides the conjectural status and asymptotic commentary
% in verification/q4-check/q4-check.md. The supplied report records completed review.

The A13 extension replaces strictly positive log-host forcing with
nonnegative host abundance. Its raw-host covariance expansion uses
abundance covariances directly, which remain defined when hosts are absent.
In the two-stage ring with one fluctuating transfer, for a stationary
ergodic process $h_t\geq0$ with $\mathbb{E}h_t=1$ and
$\mathbb{E}h_t^4<\infty$, exposed geometric mean $a>1$, and steady
geometric mean $s\geq1$, the report proves a unique finite threshold and
\begin{equation}
 D=\frac{A}{s}+O(s^{-2}),\qquad
 A=\frac{\operatorname{Var}(h_t)}{2}
   +\sum_{\ell\geq1}(1-1/a)^{\ell-1}
      \operatorname{Cov}(h_t,h_{t-\ell}).
 \label{num-zero-expansion}
\end{equation}
Here $A$ is the leading coefficient and $\ell$ is the time lag. A negative
$A$ gives eventual approach to zero from below and increasing $D$ at every
fixed multiplicative increase in sufficiently long stay. A positive $A$
gives the opposite ordering. Eventual derivative signs are proved for
bounded hosts. Their extension to all unbounded fourth-moment laws remains
open, and $A=0$ requires higher-order terms. These statements do not imply
monotonicity over every finite stay.

For stationary ergodic mean-one hosts, the same report proves the outage
upper bound for a fixed-staying ring
with one fluctuating transfer, at least one positive staying probability,
outage probability strictly between zero and one, and a finite positive
conditional geometric mean of host abundance on nonzero steps. It also
confirms the alternating-outage formula and the independence result allowing
host zeros. The bound does not generally tend to zero as one stay grows.
Its limit is $-\ln[(1-q_0)G_+]$, where $q_0$ is outage probability and
$G_+$ is the conditional positive-step geometric mean. The limit is zero
for pure binary outages.

Independent review is complete. Separate code and direct labelled-path
enumeration passed 125 of 125 principal cases, and the review accepts the
stated proofs while retaining the general derivative-sign gap. The report
does not recount the 51 empirical reversals or extend the bound to several
varying transfers. Thus the ledger's extension category is supported by
subsequent theory within explicit conditions, rather than by a claim that
every real host sequence must benefit from longer residence.

\subsubsection{Design of a biological test}
\label{num-esm-experiment}
% Source: notes/decisive-experiment-v2.md, especially sections 6--10 and Refereed.

A biological test requires independently measured stage residence and
survival distributions, encounter and disappearance hazards, offspring
production, initial composition, and observation weights. Constant-host
growth alone does not identify $R_0$ without the generation kernel.
Likewise, held-density growth measures a composite response and cannot
identify waiting versus wasted transmission without direct encounter
measurements. Manipulating density must not silently remove parasites
with hosts. If removal occurs, it belongs in the model and in controls
that receive the same handling.

The reviewed design contains five kinds of arm. A constant-mean-density
control establishes the reference. Held densities at 0.2 and 1.8 times
the mean measure growth responses alongside direct infection and loss
assays. A fast square wave with nominal period one-quarter of a generation
helps assess averaged encounter effects. A period sweep around the
measured generation kernel examines timing and regularity of residence.
Finally, random telegraph forcing (switching between two density levels
after random dwell times) examines a fading-memory environment. The
proposed nominal duration was 60 days, with counts at days 30 and 60.
These times are illustrative, pending a feasible biological system.

Telegraph schedules should be implemented in complete independent mirrored
pairs, using the same switching times with high and low phases exchanged.
Exposure is balanced over these blocks, not assumed equal in each random
path. Square-wave comparisons require prespecified phase blocks. Actual
density levels, transition ramps, and integrated exposures must be recorded.
Schedules with few realized switches should not be discarded after running,
because conditioning on switching changes the environmental law.
The period response must propagate uncertainty in the measured residence
kernels. Absence of a nominal resonance reversal is not a general rejection
of the theory or of regular stays.

The earlier design's recommendation of one to six replicates was withdrawn
after review. It counted measurement error alone and was not a power
calculation. The 51--61 per-cell estimates above instead concern a
particular simulated four-cell extinction interaction. They do not supply
replication for all five kinds of experimental arm. A usable design needs
biological and observation variation, zeros and extinction, calibration
uncertainty, complete independent blocks, and prespecified competing
mechanisms. Mean-log count contrasts, arithmetic mean counts, threshold
losses, and extinction probabilities must remain separate outcomes.

% Sources: notes/decisive-experiment-v2.md section 9; casestudies/leads/leads.md.
Candidate systems in the source review include \textit{Paramecium caudatum}
with \textit{Holospora undulata}, \textit{Pseudomonas fluorescens} SBW25
with phage $\Phi2$, and \textit{Daphnia} with \textit{Ordospora}.
The ciliate system has a described transmission stage and scheduled
culling and between-patch migration, but the review does not supply
spore survival or encounter efficiency. Phage systems permit controlled host supply, while
manipulation-associated phage losses and evolving resistance remain
unresolved for this design. The available water-flea experiments varied
temperature rather than host supply at a fixed mean. These are candidate
systems identified from literature cards and inventories, not established
experimental implementations
\citep{duncan2013,harrison2013,krichel2023}.

The later archive check identifies fixed-mean resource treatments in
\textit{Tribolium castaneum}. Constant 20\,g flour was compared with
alternating 32/8\,g, 28/12\,g, or 32/32/8/8\,g schedules in the relevant
experiments \citep{desharnais2023}. These preserve nominal mean
flour over complete cycles, not beetle abundance. None of the 11
archived datasets contains a parasite. They supply host trajectories for
the simulated world, but cannot test parasite persistence. The additional
alga--virus archive imposed dilution fluctuations only in virus-free
cultures, and the pond spore study supplies physical-settling calibration
without a manipulated fixed-mean host contrast
\citep{lepennec2024,tadiri2026}.

No system has been selected and no proposed biological experiment has
been run. The reviewed design remains conditional on a workable density
manipulation, independent calibration, an observation model, and a
replication calculation that includes their uncertainty. Completing those
steps would allow the experiment to distinguish protection from a steady
stage from changes in mean transmission, parasite removal, or development.
